\section{Resultados experimentales y análisis}
\label{sec:resultados}

% Figuras de ancho completo: se declaran al inicio para que LaTeX las ubique arriba de una página.
\begin{figure*}[t]
\centering
\includegraphics[width=\textwidth]{gradcam_ejemplos.pdf}
\caption{Grad-CAM del modelo EfficientNet-B0 con augmentation sobre imágenes de test. Las tres
columnas de la izquierda muestran aciertos típicos (imagen con la mediana de $E_{\text{hoja}}$
dentro de su clase) y las tres de la derecha, los errores con mayor confianza. El contorno
blanco indica la máscara de la hoja usada para calcular $E_{\text{hoja}}$.}
\label{fig:gradcam_ejemplos}
\end{figure*}

\subsection{Comparación de backbones (linear probe)}
\label{subsec:res_linear_probe}
La Tabla~\ref{tab:linear_probe} resume el desempeño de una capa lineal entrenada sobre los
embeddings congelados de cada backbone. Sin modificar ningún peso preentrenado, todas las
arquitecturas superan el 94\,\% de macro-F1 en test, lo que confirma que las representaciones
de ImageNet transfieren bien a este dominio. Las arquitecturas diseñadas para dispositivos
móviles superan a ResNet18 (macro-F1 0,960 contra 0,942) a pesar de tener entre 4 y 12 veces
menos parámetros. Entre ellas, las diferencias en test son menores a 0,002, del orden de la
variabilidad entre validación y test: EfficientNet-B0 obtiene el mayor macro-F1 en validación
y MobileNetV3-Small, uno prácticamente igual con una extracción 3,3 veces más rápida. La
entropía cruzada ponderada no mostró un beneficio consistente: modificó el macro-F1 en test
entre $-0{,}007$ y $+0{,}004$ según el modelo, por lo que no la utilizamos en el fine-tuning.

\begin{table}[t]
\centering
\caption{Linear probe sobre backbones congelados (entropía cruzada sin ponderar)}
\label{tab:linear_probe}
\setlength{\tabcolsep}{4pt}
\begin{tabular}{@{}lccccc@{}}
\toprule
\textbf{Backbone} & \textbf{Parám.} & \textbf{Extracción} & \multicolumn{2}{c}{\textbf{Macro-F1}} & \textbf{Bal. acc.} \\
\cmidrule(lr){4-5}
 & (M) & (img/s) & val & test & test \\
\midrule
MobileNetV3-Small & 0,9  & 208,6 & 0,972          & 0,960          & 0,962 \\
MobileNetV3-Large & 3,0  & 91,1  & 0,967          & 0,960          & 0,957 \\
EfficientNet-B0   & 4,0  & 62,7  & \textbf{0,974} & 0,959          & 0,961 \\
ResNet18          & 11,2 & 55,4  & 0,943          & 0,942          & 0,941 \\
\bottomrule
\end{tabular}
\end{table}

\subsection{Fine-tuning con y sin data augmentation}
\label{subsec:res_finetuning}
El fine-tuning parcial mejora de forma sustancial el linear probe
(Tabla~\ref{tab:finetuning}). El macro-F1 en test pasa de 0,960 a 0,986 en MobileNetV3-Small
y de 0,959 a 0,994 en EfficientNet-B0, lo que equivale a reducir la fracción de error
($1-\text{macro-F1}$) un 65\,\% y un 85\,\%, respectivamente. Adaptar las últimas etapas del
backbone resulta, entonces, mucho más determinante que la elección de la arquitectura en la
etapa de linear probe. EfficientNet-B0 supera a MobileNetV3-Small en todas las métricas, con
la mitad de errores, pero su entrenamiento demora tres veces más.

El data augmentation no mejoró el desempeño en el test de PlantVillage: las diferencias
entre las corridas con y sin augmentation son de a lo sumo 0,002 en macro-F1 (6 imágenes de
2\,949), un margen que no consideramos significativo con una sola semilla. Sí se observa su
efecto regularizador: sin augmentation, EfficientNet-B0 termina con una pérdida de
entrenamiento de 0,003 frente a 0,016 en validación (memoriza el conjunto de entrenamiento),
mientras que con augmentation la pérdida de entrenamiento (0,038) queda por encima de la de
validación (0,019). Como el test proviene de la misma distribución que el entrenamiento, es
esperable que el beneficio del augmentation aparezca recién ante condiciones de captura
distintas (Secciones~\ref{subsec:res_robustez} y~\ref{subsec:res_externos}). Por último, en
las cuatro corridas el macro-F1 de validación todavía crecía en la última época, por lo que
un entrenamiento más largo podría mejorar levemente estos valores.

\begin{table}[t]
\centering
\caption{Fine-tuning parcial a $160\times160$ (8 épocas, CPU)}
\label{tab:finetuning}
\setlength{\tabcolsep}{3.5pt}
\begin{tabular}{@{}lcccccc@{}}
\toprule
\textbf{Modelo} & \textbf{Aug.} & \textbf{Parám.} & \textbf{F1 val} & \textbf{F1 test} & \textbf{Errores} & \textbf{Tiempo} \\
 & & entren. (M) & & & test & (min) \\
\midrule
MobileNetV3-S   & no & 1,39 & 0,990          & 0,986          & 24          & 10,7 \\
MobileNetV3-S   & sí & 1,39 & 0,987          & 0,986          & 30          & 10,6 \\
EfficientNet-B0 & no & 3,17 & 0,995          & \textbf{0,994} & \textbf{11} & 33,8 \\
EfficientNet-B0 & sí & 3,17 & \textbf{0,995} & 0,993          & 17          & 32,4 \\
\bottomrule
\end{tabular}
\end{table}

\subsection{Análisis por clase y de errores}
\label{subsec:res_errores}
Con el modelo seleccionado (EfficientNet-B0 con augmentation, el de mayor macro-F1 en
validación), todas las clases superan un F1 de 0,975. Las más bajas son tizón temprano del
tomate (0,975) y papa sana (0,976). Esta última es la clase con menor F1 en tres de las
cuatro corridas, lo que es coherente con que sea la clase minoritaria (21 imágenes de test).
De los 17 errores, 16 ocurren entre enfermedades de una misma especie y solo uno confunde una
hoja sana con una enferma. Las confusiones más frecuentes se dan entre enfermedades del
tomate con síntomas visualmente parecidos: tizón temprano contra tizón tardío, y mancha diana
contra ácaro araña o moho foliar. Estos mismos pares ya aparecían como los principales
errores del linear probe, lo que sugiere una ambigüedad propia de los síntomas más que una
limitación de un modelo en particular.

\subsection{Explicabilidad (Grad-CAM)}
\label{subsec:res_gradcam}
% Nota: Grad-CAM se calculó sobre un reentrenamiento de EfficientNet-B0 con augmentation
% (branch gradcam). Actualizar al unificar los resultados en main.
El análisis de Grad-CAM se realizó sobre un reentrenamiento del mismo modelo (macro-F1 en
test de 0,994; 13 errores). En el conjunto de test completo, el 67,8\,\% de la activación
recae sobre la hoja, que ocupa en promedio el 40,2\,\% de la imagen: la atención está
1,77 veces más concentrada en la hoja de lo que se esperaría por azar. Además, el máximo de
activación cae dentro de la hoja en el 97,0\,\% de las imágenes (Fig.~\ref{fig:energia_hoja}).
Los aciertos de la Fig.~\ref{fig:gradcam_ejemplos} muestran que el modelo se enfoca en las
lesiones necróticas y en las zonas cloróticas.

\begin{figure}[t]
\centering
\includegraphics[width=\columnwidth]{gradcam_energia_hoja.pdf}
\caption{Distribución de $E_{\text{hoja}}$ (fracción de la activación Grad-CAM sobre la hoja)
en las 2\,949 imágenes de test. Los triángulos marcan los 13 errores; la línea punteada, el
área media que ocupa la hoja.}
\label{fig:energia_hoja}
\end{figure}

Los errores no se explican por una atención desviada hacia el fondo: su $E_{\text{hoja}}$
medio (0,63) y su \emph{pointing game} (92\,\%) son apenas menores que los de los aciertos
(0,68 y 97\,\%). Es decir, el modelo mira la hoja pero interpreta mal síntomas ambiguos. La
excepción es una hoja de TYLCV de tamaño muy reducido, en la que la activación se extiende
sobre el fondo (Fig.~\ref{fig:gradcam_ejemplos}, cuarta columna). Las clases con menor
$E_{\text{hoja}}$ son el virus del mosaico (0,54) y el tizón tardío del tomate (0,55), cuyas
hojas suelen estar enrolladas o deformadas. Parte de la activación fuera de la hoja se debe a
la baja resolución del mapa ($5\times5$), que al interpolarse se extiende sobre los bordes.
Aun así, que cerca de un tercio de la activación caiga fuera de la hoja indica que el fondo
no es completamente irrelevante para el modelo, en línea con el sesgo descrito
en~\cite{noyan2022uncovering}.

\subsection{Robustez ante perturbaciones}
\label{subsec:res_robustez}

\subsection{Evaluación con datos externos}
\label{subsec:res_externos}

\subsection{Costo computacional}
\label{subsec:res_costo}
La Tabla~\ref{tab:costo} compara el costo de inferencia de los dos modelos ajustados en una
CPU de notebook (Intel Core i7-1355U), con entrada de $160\times160$. Ambos permiten
clasificar en tiempo real sin GPU: EfficientNet-B0 procesa una imagen en 18\,ms y
MobileNetV3-Small en 7,7\,ms, con un modelo 2,6 veces más liviano (5,9\,MB). Si se prioriza
el desempeño, EfficientNet-B0 es la mejor opción para el prototipo. MobileNetV3-Small, que
comete el doble de errores pero es 2,3 veces más rápido, resulta una alternativa adecuada
para dispositivos más limitados, como un teléfono de gama baja.

\begin{table}[t]
\centering
\caption{Costo de inferencia en CPU ($160\times160$, fp32)}
\label{tab:costo}
\setlength{\tabcolsep}{3pt}
\begin{tabular}{@{}lcccc@{}}
\toprule
\textbf{Modelo} & \textbf{Parám.} & \textbf{Tamaño} & \textbf{Latencia bs=1} & \textbf{Throughput} \\
 & (M) & (MB) & media (p95) [ms] & bs=32 [img/s] \\
\midrule
MobileNetV3-Small & 1,53 & 5,9  & 7,7 (9,1)   & 920 \\
EfficientNet-B0   & 4,03 & 15,5 & 18,0 (25,8) & 152 \\
\bottomrule
\end{tabular}
\end{table}
